\documentclass[11pt]{article}
\usepackage{amsmath,amssymb,amsthm,mathtools,hyperref,enumitem,geometry}
\geometry{margin=1.1in}
\usepackage[T1]{fontenc}
\usepackage{microtype}
\usepackage{mathrsfs}
\usepackage{listings}
\lstset{basicstyle=\ttfamily\small,breaklines=true}

\title{Computing the Hodge Lattice of Fermat Hypersurfaces:\\
Rank and Discriminant via Shioda--Ran and Smith Normal Form}
\author{Prepared for a friendly duel with a persistent proposer}
\date{\today}

\newtheorem{theorem}{Theorem}
\newtheorem{lemma}{Lemma}
\newtheorem{prop}{Proposition}
\newtheorem{cor}{Corollary}
\theoremstyle{definition}
\newtheorem{definition}{Definition}
\newtheorem{remark}{Remark}

\newcommand{\Xmn}{X_m^n}
\newcommand{\PP}{\mathbf{P}}
\newcommand{\ZZ}{\mathbf{Z}}
\newcommand{\QQ}{\mathbf{Q}}
\newcommand{\CC}{\mathbf{C}}
\newcommand{\muone}{\mu_m^{\,n+2}/\mu_m}
\newcommand{\Hdg}{\mathrm{Hdg}}
\newcommand{\NS}{\mathrm{NS}}
\newcommand{\rank}{\mathrm{rank}}
\newcommand{\disc}{\mathrm{disc}}

\begin{document}
\maketitle

\begin{abstract}
We give a practical, computer-ready blueprint for computing the \emph{rank} and, when feasible, the \emph{discriminant} of the lattice of Hodge cycles in the middle cohomology of Fermat hypersurfaces
\(\Xmn : x_0^m+\cdots+x_{n+1}^m=0\subset \PP^{n+1}\) with \(n\) even.
Two complementary routes are described:
\begin{itemize}[nosep]
\item[(A)] a purely combinatorial count (after Shioda--Ran) giving \(\mathrm{rank}\) as \(1+\#\mathfrak{B}_m^{\,n}\);
\item[(B)] a Smith normal form (SNF) computation of a Gram matrix built either from (B1) linear algebraic cycles \( \PP^p \subset \Xmn\), or (B2) vanishing cycles projected to the Hodge subspace by character idempotents.
\end{itemize}
We also package ready-to-run \textsf{SageMath} code that implements (A) in full generality and (B1) for fourfolds. Worked examples include the Fermat surface of degree \(13\) and Fermat fourfolds of degrees \(5,7,9\).
\end{abstract}

\section{Set-up and the Shioda--Ran description}
Let \(n=2p\) be even and \(X=\Xmn\). Write \(G=\mu_m^{\,n+2}/\mu_m\) for the diagonal symmetry and \(\widehat G\cong \{(a_0,\dots,a_{n+1})\in (\ZZ/m\ZZ)^{n+2}\mid \sum a_i\equiv 0\}\) for its character group.
Following Shioda and Ran, the primitive middle cohomology decomposes into one-dimensional character eigenspaces \(V(\alpha)\) and the primitive Hodge classes are a direct sum of those \(V(\alpha)\) picked out by a simple congruence condition:

\begin{definition}[Shioda's set \(\mathfrak U_m^{\,n}\) and \(\mathfrak B_m^{\,n}\)]
Let \(\langle a\rangle\) denote the representative of \(a\in\ZZ/m\ZZ\) in \(\{1,\dots,m-1\}\). Put
\begin{align*}
\mathfrak U_m^{\,n}&=\Big\{\alpha=(a_0,\dots,a_{n+1})\in\widehat G\;:\; a_i\not\equiv 0\ \forall i\Big\},\\
\mathfrak B_m^{\,n}&=\Big\{\alpha\in\mathfrak U_m^{\,n}\;:\; \sum_{i=0}^{n+1}\big\langle t a_i\big\rangle =(p+1)m\ \text{ for all }t\in (\ZZ/m\ZZ)^\times\Big\}.
\end{align*}
\end{definition}

\begin{theorem}[Ran--Shioda]
If \(n=2p\) then the primitive Hodge classes are
\(\Hdg^p_{\mathrm{prim}}(X)=\bigoplus_{\alpha\in \mathfrak B_m^{\,n}} V(\alpha)\), each summand one-dimensional over \(\CC\). Consequently,
\begin{equation}\label{eq:rank-formula}
\rank\,\Hdg^p(X;\ZZ)=1+\#\,\mathfrak B_m^{\,n},
\end{equation}
the extra \(1\) being the non-primitive class \(h^p\) of the hyperplane.\footnote{See, e.g., da~Silva's exposition \emph{Notes on the Hodge conjecture for Fermat varieties}, Theorem~2.1 and \S2 \cite{daSilva2021}, which restates Shioda's original result \cite{Shioda1979} and Ran \cite{Ran1980}.}
\end{theorem}

\paragraph{Algorithm A (rank).} The condition defining \(\mathfrak B_m^{\,n}\) can be enforced with a small integer linear system on the \emph{multiplicity vector} \((c_1,\dots,c_{m-1})\), where \(c_r\) counts how many coordinates equal \(r\). Writing \(k=n+2\) and \(U=(\ZZ/m\ZZ)^\times\), one requires
\begin{equation}\label{eq:ILP}
\sum_{r=1}^{m-1} c_r = k,\qquad \sum_{r=1}^{m-1}\big\langle t r\big\rangle c_r = (p+1)m\quad(\forall\,t\in U),\qquad c_r\in\ZZ_{\ge 0}.
\end{equation}
For each solution \(c\), the number of ordered \((a_0,\dots,a_{k-1})\) it represents is the multinomial \(\binom{k}{c_1,\dots,c_{m-1}}\). Summing these multinomials over all solutions of \eqref{eq:ILP} yields \(\#\,\mathfrak B_m^{\,n}\), hence the rank by \eqref{eq:rank-formula}. This is implemented in the attached \textsf{Sage} file (function \texttt{hodge\_rank\_via\_B}).

\medskip
\section{Intersection of linear cycles and a Gram matrix}
For \(n=2p\), the Fermat equation splits coordinates into \(p+1\) pairs; fixing a pairing and \(\zeta_j\in\mu_m\) per pair, the linear subspace
\begin{equation}\label{eq:linear-cycle}
L(\boldsymbol\zeta):\quad x_{2j}=\zeta_j\,x_{2j+1}\qquad (j=0,\dots,p)
\end{equation}
lies on \(\Xmn\). Let \(P_i\) range over all such \(L(\boldsymbol\zeta)\) for all pairings and choices of \(\boldsymbol\zeta\) (there are \( \frac{(n+2)!}{2^{p+1}(p+1)!}\cdot m^{p+1}\) of them). The intersection numbers inside \(X\) depend only on \(d=m\) and on the \emph{projective} dimension \(r=\dim(P_i\cap P_j)\) of their intersection \emph{in} \(\PP^{n+1}\):
\begin{equation}\label{eq:AMV-formula}
P_i\cdot P_j=\frac{1-(-d+1)^{\,r+1}}{d}\qquad(r=-1,0,1,\dots,p),
\end{equation}
with the convention \((\cdot)=0\) for $r=-1$. This closed formula (no Gr\"obner bases needed) is stated in \cite[Eq.\,(6)]{AMV2019}.

\paragraph{Primitive part.} Following \cite[\S3.1]{AMV2019}, pass to primitive classes \( \check P_i:=P_i-P_1\) and form the integer Gram matrix \(A_1=[\check P_i\cdot \check P_j]\). Its SNF gives the elementary divisors of the lattice generated by the linear cycles in primitive middle homology.

\paragraph{From primitive to full Hodge lattice.}
As explained in \cite[\S3.3]{AMV2019}, to pass from the primitive lattice to the full Hodge lattice one \emph{decreases} the multiplicity of the elementary divisor \(1\) by \(2\) and \emph{adds} a new elementary divisor \(d=m\). (Equivalently, adjoin \(h^p\) and apply SNF.)

\paragraph{Primitivity and discriminant.}
In favorable cases (e.g.\ degree \(5\) fourfolds, and many Fermat surfaces), one can certify that the sublattice generated by linear cycles is \emph{primitive} in the Hodge lattice using supersingular reduction and the Artin--Tate apparatus; see \S6 below and \cite{AMV2019,SchuettShiodaVanLuijk2009}. When primitivity holds, the SNF of \(A_1\) (after the preceding adjustment) yields the elementary divisors and hence the \emph{discriminant}.

\section{Alternative route: vanishing cycles and idempotent projectors}
One can compute the Hodge lattice intrinsically, avoiding algebraic cycles, by starting with a Pham basis of vanishing cycles for the Brieskorn--Pham singularity \(x_0^m+\cdots+x_{n+1}^m\), letting \(G\) act, and projecting to the Hodge subspace by
\[
E=\sum_{\alpha\in\mathfrak B_m^{\,n}} e_\alpha,\qquad
e_\alpha=\frac{1}{|G|}\sum_{g\in G}\chi_\alpha(g)^{-1} g.
\]
Restricting the vanishing-cycle intersection form to \(\mathrm{im}(E)\) gives a Gram matrix for the primitive Hodge lattice. This is the second half of \cite{AMV2019} (their Algorithm~2), where the authors compare SNFs from the two routes (linear vs.\ vanishing cycles) to prove equality of lattices.

\section{Worked examples}
\subsection*{Surface, degree \(13\)}
The Fermat surface \(S: x_0^{13}+x_1^{13}+x_2^{13}+x_3^{13}=0\) has \(\rho(S)=397\) and \(\disc \NS(S)=13^{300}\). This is proved in \cite[\S6.5]{SchuettShiodaVanLuijk2009} using supersingular reduction at \(p=5\) and the Artin--Tate formula (see the PDF for the exact statement and numbers).

\subsection*{Fourfold, degree \(5\)}
For \(X_5^4\) one has \(\rank \Hdg^2(X_5^4)=1+\#\mathfrak B_5^4=401\) (Algorithm~A). Aljovin--Movasati--Villaflor compute the SNF of the Hodge lattice and list its elementary divisors as
\[
1^{166}\cdot 5^{174}\cdot 25^{54}\cdot 125^{7},
\]
so the discriminant is \(5^{174+2\cdot 54+3\cdot 7}=5^{303}\) \cite[Table~1, row $(n,d)=(4,5)$]{AMV2019}.

\subsection*{Fourfolds, degrees \(7\) and \(9\)}
Algorithm~A gives
\[
\#\mathfrak B_7^4=1860,\qquad \#\mathfrak B_9^4=8000,
\]
hence \(\rank \Hdg^2(X_7^4)=1861\) and \(\rank \Hdg^2(X_9^4)=8001\). The attached \textsf{Sage} code produces these counts in milliseconds. For discriminants, one can proceed via \S3 using all linear \(\PP^2\)'s; primitivity (and thus equality with the Hodge lattice) would require an additional argument as in \cite{AMV2019}.

\section{What the attached \textsf{SageMath} file does}
The file \texttt{fermat\_hodge\_lattice.sage} contains:
\begin{itemize}
\item \texttt{hodge\_rank\_via\_B(m,n)}: full implementation of Algorithm~A, counting \(\#\mathfrak B_m^{\,n}\) by solving \eqref{eq:ILP} via a tiny exhaustive search on multiplicities \(c_r\). For $(n,d)=(4,5),(4,7),(4,9)$ it returns \(400,1860,8000\).
\item \texttt{linear\_cycles\_intersection\_number(c1,c2,m)}: computes \(r=\dim(P_i\cap P_j)\) by linear algebra over \(\QQ(\zeta_m)\) and returns the integer \( (1-(-m+1)^{r+1})/m\) as in \eqref{eq:AMV-formula}.
\item \texttt{linear\_cycles\_SNF(m)} (fourfolds): builds a spanning set of linear \(\PP^2\)'s, forms the Gram matrix of the primitive classes, extracts a full-rank submatrix, and takes its SNF. A helper \texttt{adjust\_primitive\_to\_full} performs the \(+\,h^2\) adjustment explained above.
\end{itemize}
The code is written to be readable and modifiable; it uses only standard \textsf{Sage} functionality (cyclotomic fields, integer matrices, SNF).

\paragraph{Practical tip.} For \(m=5\) the full set of linear cycles has size \(15\cdot 5^3=1875\). You do \emph{not} need them all: the script greedily selects cycles until the target rank is reached (401), then computes SNF on that smaller matrix. For \(m=7,9\) the same strategy is advisable.

\bigskip
\noindent\textbf{Acknowledgements.} The algorithms and the idea of comparing SNFs are due to Aljovin, Movasati, and Villaflor; Shioda and Ran underlie the rank computation and the representation-theoretic picture. Any implementation errors in the accompanying code are mine.

\begin{thebibliography}{99}

\bibitem{Shioda1979}
T.~Shioda, \emph{The Hodge Conjecture for Fermat Varieties}, Math.\ Ann.\ \textbf{245} (1979), 175--184.

\bibitem{Ran1980}
Z.~Ran, \emph{Cycles on Fermat hypersurfaces}, Compositio Math.\ \textbf{42} (1980), 121--142.

\bibitem{AMV2019}
E.~Aljovin, H.~Movasati, R.~Villaflor~Loyola, \emph{Integral Hodge conjecture for Fermat varieties}, arXiv:1711.02628v2 (2019) and Exp.\ Math.\ (2019). Algorithm boxes and Tables~1--2 give the SNF workflows and the elementary divisors for quartic/quintic fourfolds.

\bibitem{SchuettShiodaVanLuijk2009}
M.~Sch\"utt, T.~Shioda, R.~van~Luijk, \emph{Lines on Fermat surfaces} (2009 preprint). Section~6.5 proves for $m=13$ that $\rho=397$ and $\disc \NS=13^{300}$.

\bibitem{daSilva2021}
G.~da~Silva~Jr, \emph{Notes on the Hodge conjecture for Fermat varieties}, Experimental Results (2021). Theorem~2.1 reviews the $\Hdg$-decomposition via the set $\mathfrak B_m^{\,n}$.

\end{thebibliography}

\end{document}